\subsection{The directional germ and a single unknown launch law}
The principal inverse uses one stationary launch density $j$, with
unknown compact convex support $A$. The controller prescribes a nominal
center $x$, a direction $n_\theta$ and a length $t>0$, and observes only
the bit $B_{tn_\theta}(x+Z)$, where $Z$ has density $j$. The exact datum
is the family of raw forward means $F_{t,\theta}(x)$ for all centers,
directions and arbitrarily short positive lengths. Directions are
command labels. No launch realization or impact record is observed.

For locally finite separated $C^2$ strictly convex obstacles, the
first-order germ exists in $L^1_{\mathrm{loc}}$:
\begin{equation}\label{eq:intro-single-germ}
 q_\theta=\lim_{t\downarrow0}\frac{F_{t,\theta}}t.
\end{equation}
Let $c_C(\phi)$ and $r_C(\phi)$ be the boundary point with normal
$n_\phi$ and its radius of curvature. The collision-strip formula and
the planar angular identity give
\begin{equation}\label{eq:intro-single-normal}
 (\partial_\theta^2+1)q_\theta(x)
 =\sum_{C,\,\sigma=\pm1}r_C(\theta+\sigma\pi/2)
                     j(c_C(\theta+\sigma\pi/2)-x).
\end{equation}
The right side is a positive measure density. Its spatial components
resolve the antipodal pair left by angular inversion. When the
footprint diameter is smaller than the inter-obstacle separation and
every obstacle width, each component is exactly one translate of
$-A$. Its normalized density recovers $j$ in the centered footprint
frame. The union of its Steiner points over directions recovers the
entire obstacle boundary. Thus
Theorem~\ref{thm:single-law-rigidity} identifies the canonical triple
\begin{equation}\label{eq:intro-single-triple}
       (\mathcal O-s(A),\ A-s(A),\ j(\,\cdot+s(A))).
\end{equation}
It uses neither a second launch setting nor a homothety relation.
The complete observational fiber is the common translation of the
table and density, as stated in Corollary~\ref{cor:single-law-fiber}.
The first-order germ consequently determines every finite-length
forward response and the full translation period group.

Theorem~\ref{thm:single-resolved-obstacle} gives a broader exact
form: it suffices that $\Delta<d$ and one obstacle have diameter
larger than $\Delta$. That obstacle need not be identified in
advance, and no lower width bound is required for the others.
A minimum-area angular support component recovers the common density;
the odd part of the germ recovers occupation and hence the remaining
obstacles. The finite theorem below retains uniform separation of
all angular copies.

The exact datum in \eqref{eq:intro-single-germ} remains a field on the
whole plane. Its limit and angular derivatives are operations on exact
means. Section~\ref{sec:single-law-finite} gives a separate finite
experiment that estimates regularized versions of
\eqref{eq:intro-single-normal} from the original attempted bits.
Spatial and angular smoothing produce a positive comparison field;
integration by parts places all derivatives on known kernels. The
finite-length bias is bounded weakly, without a modulus of continuity
or an upper bound for the unknown density.
For $s=6+\beta$ and the boundary-mass exponent $\gamma$, its
explicit sufficient attempted-bit bound is
\begin{equation}\label{eq:intro-single-finite-cost}
 N_\nu\le C\nu^{-Q_\gamma}\log\frac{C}{\nu\delta},
 \qquad Q_\gamma=\frac{(3\gamma+27/2)s}{s-2}.
\end{equation}
This rate implements the joint unknown-footprint reconstruction;
its necessity is not asserted.
 The small probability of a
collision after spatial averaging gives a variance bound and an
explicit Bernstein sampling cost. A finite reconstruction of both
the footprint and obstacle boundaries follows under the quantitative
smoothness, boundary-mass and aperture assumptions stated there.

\subsection{The retained pooled experiment and its sharp power}
The reciprocal compass experiment uses a different, more restricted
command family. Its difference field $g=F-R$ determines occupation
through a stopped-Poisson inverse. The fixed-stencil formulation in
Theorem~\ref{thm:finite-stencil} is pointwise: its site count depends
on $(D+\Delta)/t$, and costs for several targets are added. It neither
supplies the directional derivatives in \eqref{eq:intro-single-normal}
nor makes a finite global periodicity assertion.

Section~\ref{sec:sharp-stationary} retains the sharp stationary power
for one fixed known uniform-disk law. Put
\[
                    q_0=\frac{3s/2+1}{s-2},\qquad s=6+\beta.
\]
On the same nondegenerate bounded physical class, at the same fixed
confidence and with worst-case expected attempts as cost,
Corollary~\ref{cor:stationary-minimax} gives
\begin{equation}\label{eq:intro-single-benchmark}
 c\nu^{-q_0}\le N^*_{\rm short}(\nu,\delta)
       \le N^*_{\rm pool}(\nu,\delta)
       \le C\nu^{-q_0}\log\frac{C}{\nu\delta}.
\end{equation}
The lower design permits adaptive directions and lengths tending to
zero. The upper design uses a fixed pooled compass with
$2t+2R_+<d_0$; its deterministic attempt cap also bounds expected
cost. One logarithmic factor remains. This fixed-known-law result is
not a minimax assertion for the new joint unknown-density inverse.

\subsection{Data, finite resources and organization}
Each attempted preparation is charged, including a solid start or a
miss. We write $N$ for attempted bits, $J$ for command-labelled batches
with revisits counted again, and $S$ for distinct nominal spatial
sites. Numerical description length, arithmetic and physical
positioning are separate resources. In particular, random centers
used to estimate a scalar integral are counted as well as boundary
search centers.

The new germ inverse uses raw forward means. In the retained
constructions, occupation and the area normalization use reciprocal
differences, whereas width and perimeter deficits use a raw forward
mean. The homothetic inverses continue to assume labelled homothetic
settings and recover their ratios and registrations. Their proofs
remain complete, but that apparatus relation is not a hypothesis of
the principal single-law theorem.

Exact period equality is a whole-field statement without a periodicity
prior. Uniform finite primitive-lattice decisions use the bounded
periodic class and its positive patch margin. A finite nonperiodic
configuration can instead be recovered in a protected aperture known
to contain it completely. These acquisition assumptions appear in
the finite statements; none asserts a finite test for crystallinity
of an arbitrary unknown infinite configuration.

Sections~\ref{sec:single-law-rigidity} and
\ref{sec:single-law-finite} contain the principal exact inverse and
its finite realization. Section~\ref{sec:sharp-stationary} provides
the fixed-known-law benchmark, and Section~\ref{sec:comparison}
places the methods relative to earlier work. The appendices contain
all auxiliary and retained constructions: global and finite-stencil
occupation, stationary boundary queries, homothetic calibration,
period recognition, physical packing and finite controls. The
directional-germ proof does not depend on the stopped-Poisson inverse.
